\documentclass[11pt]{article}

\usepackage[margin=1in]{geometry}
\usepackage{amsmath, amssymb}
\usepackage{xcolor}
\usepackage{listings}
\usepackage{hyperref}
\usepackage{array}
\usepackage{booktabs}

\hypersetup{
    colorlinks=true,
    linkcolor=blue!60!black,
    urlcolor=blue!60!black,
    citecolor=blue!60!black
}

\lstdefinestyle{bugpython}{
  language=Python,
  basicstyle=\ttfamily\small,
  keywordstyle=\color{blue!60!black},
  stringstyle=\color{green!40!black},
  commentstyle=\color{gray!70!black},
  showstringspaces=false,
  breaklines=true,
  frame=single,
  rulecolor=\color{gray!30},
  columns=fullflexible,
  keepspaces=true
}

\title{SymPy Bug Report}
\author{}
\date{}

\begin{document}

\maketitle

\section{Bug 23: solveset returns 0 for atan(x) = pi outside the principal atan range}

\subsection{Minimal Reproducer}

\medskip

\begin{lstlisting}[style=bugpython]
from sympy import Eq, S, N, atan, pi, solveset, symbols

x = symbols("x")
eq = Eq(atan(x), pi)
sol = solveset(eq, x, domain=S.Complexes)

print("solution =", sol)
for s in sol:
    print("residual at", s, "=", N(atan(s) - pi, 80))
\end{lstlisting}

\subsection{Observed Behavior}

\medskip

\begin{lstlisting}[style=bugpython]
solution = {0}
residual at 0 = -3.1415926535897932384626433832795028841971693993751058209749445923078164062862090
\end{lstlisting}

SymPy returns the candidate \(x = 0\).  Substituting that value into the
original equation gives
\[
  \operatorname{atan}(0) - \pi = -\pi,
\]
so the returned value does not satisfy the equation.

\subsection{Expected Behavior}

The correct solution set over \(S.\mathrm{Complexes}\) is \(\varnothing\).  The
principal value of \(\operatorname{atan}(0)\) is \(0\), not \(\pi\), and the
principal branch of \(\operatorname{atan}\) does not attain the real value
\(\pi\).

\subsection{Mathematical Explanation}

The equation being solved is
\[
  \operatorname{atan}(x) = \pi.
\]
Applying \(\tan\) to both sides gives the tempting candidate
\[
  x = \tan(\pi) = 0.
\]
That inference is invalid for the principal inverse tangent.  The identity
\(\tan(\operatorname{atan}(x)) = x\) holds where the principal arctangent is
defined, but \(\operatorname{atan}(\tan(y))\) does not equal \(y\) for arbitrary
\(y\).  Instead, \(\operatorname{atan}\) returns a principal value with real part
in the principal arctangent strip.  In particular,
\[
  \operatorname{atan}(0) = 0,
\]
so \(x = 0\) satisfies \(\operatorname{atan}(x) = 0\), not
\(\operatorname{atan}(x) = \pi\).  The reported answer is therefore an
extraneous solution introduced by treating a principal-branch inverse as though
it were a globally valid inverse.

\subsection{Source-Level Diagnosis}

\begin{sloppypar}
The diagnosis identifies the relevant solver code in
\nolinkurl{sympy/solvers/solveset.py}, specifically \texttt{\_solveset} and
\texttt{\_invert\_complex}.  The traced call path starts from the complex-domain
call to solve \(\operatorname{atan}(x) = \pi\).  The equation is normalized to
\texttt{atan(x) - pi}.  This expression is not handled by \texttt{\_solve\_trig},
so \texttt{\_solveset} reaches the generic inversion call.  The inverter removes
the additive \(-\pi\), obtains \texttt{atan(x) = pi}, and dispatches to the
generic inverse-function branch in \texttt{\_invert\_complex}.
\end{sloppypar}

The key rule in \(\texttt{\_invert\_complex}\), located at
\nolinkurl{sympy/solvers/solveset.py:550-557}, is:

\begin{lstlisting}[style=bugpython]
if hasattr(f, 'inverse') and f.inverse() is not None and \
   not isinstance(f, TrigonometricFunction) and \
   not isinstance(f, HyperbolicFunction) and \
   not isinstance(f, exp):
    ...
    return _invert_complex(f.args[0],
                           imageset(Lambda(n, f.inverse()(n)), g_ys), symbol)
\end{lstlisting}

For \(f = \operatorname{atan}(x)\), the inverse method returns \(\tan\), with the
inverse defined in \nolinkurl{sympy/functions/elementary/trigonometric.py:2775}.
The diagnosis notes that inverse trigonometric functions are not excluded by
this branch.  With target set \(\{\pi\}\), the branch maps the target through
\(\tan\), producing \(\{\tan(\pi)\} = \{0\}\).  The finite candidate is then
accepted because the original expression is defined at \(x = 0\); no residual
equality check rejects it.  Thus the solver returns a solution of
\(\tan(\operatorname{atan}(x)) = \tan(\pi)\), not a solution of the original
principal-value equation.

A plausible fix direction supported by the diagnosis is to treat inverse
trigonometric functions as partial inverses.  The solver could invert them only
when the target is known to lie in the appropriate principal range, or it could
post-filter finite candidates against the original equation.

\subsection{Additional Failing Instances}

The same mechanism appears for other real multiples of \(\pi\) outside the
principal range of \(\operatorname{atan}\).  Applying \(\tan\) to each target
produces \(0\), but the principal value \(\operatorname{atan}(0)\) is \(0\), not
the requested nonzero multiple of \(\pi\).  The independent verification records
the following parametrized cases.

\begin{center}
\small
\renewcommand{\arraystretch}{1.3}
\begin{tabular}{@{}>{\raggedright\arraybackslash}p{0.44\linewidth}>{\raggedright\arraybackslash}p{0.23\linewidth}>{\raggedright\arraybackslash}p{0.23\linewidth}@{}}
\toprule
\textbf{Input / Expression} & \textbf{SymPy behavior} & \textbf{Expected behavior} \\
\midrule
\(\operatorname{atan}(x) = \pi\) & returns a set containing \(0\) & \(\varnothing\) \\
\(\operatorname{atan}(x) = -\pi\) & returns a set containing \(0\) & \(\varnothing\) \\
\(\operatorname{atan}(x) = 2\pi\) & returns a set containing \(0\) & \(\varnothing\) \\
\(\operatorname{atan}(x) = -2\pi\) & returns a set containing \(0\) & \(\varnothing\) \\
\(\operatorname{atan}(x) = 3\pi\) & returns a set containing \(0\) & \(\varnothing\) \\
\(\operatorname{atan}(x) = -3\pi\) & returns a set containing \(0\) & \(\varnothing\) \\
\bottomrule
\end{tabular}
\end{center}

The expected behavior is the same across the family because each target is a
real value outside the real principal range of \(\operatorname{atan}\), while the
returned candidate is only a solution of the tangent equation obtained after an
invalid global inversion step.

\subsection{Related Issues Assessment}

The available deduplication evidence found related background on principal
branches of inverse trigonometric functions and the SymPy \texttt{solveset}
contract, but no public report of this exact reproducer or output.  The verdict
is therefore a new instance of a known failure mode: the broad issue of
respecting principal branches is known, but this exact \(\operatorname{atan}\)
case with the false result \(\{0\}\) appears previously unreported and remains
individually actionable as a focused issue and regression test.

\subsection{Proposed SymPy Regression Test}

\medskip

\begin{lstlisting}[style=bugpython]
import pytest
from sympy import Eq, S, atan, pi, solveset, symbols, simplify

x = symbols("x")

CASES = [
    (pi, 0),
    (-pi, 0),
    (2*pi, 0),
    (-2*pi, 0),
    (3*pi, 0),
    (-3*pi, 0),
]


@pytest.mark.parametrize("target,bad", CASES)
def test_solveset_atan_rejects_values_outside_principal_range(target, bad):
    sol = solveset(Eq(atan(x), target), x, domain=S.Complexes)

    assert sol == S.EmptySet
    assert simplify(atan(bad) - target) != 0
\end{lstlisting}

\end{document}
